\documentclass[12pt,a4paper,oneside]{article}

\usepackage[QX]{polski}

\usepackage[utf8]{inputenc}
\usepackage{latexsym}
\usepackage{tgpagella}
\usepackage{lmodern}
\usepackage{amsmath,amsthm,amsfonts,amssymb,alltt}
\usepackage{epsfig}
\usepackage{pdflscape}
\usepackage{caption}
\usepackage{indentfirst}
\usepackage{float}
%\usepackage{showkeys}
\bibliographystyle{plabbrv}
\usepackage{multirow}

\usepackage{color}
\usepackage[polish]{babel}
\usepackage{datetime2}
\usepackage[x11names,dvipsnames,table]{xcolor}
\usepackage{hyperref}
\hypersetup{
pdfauthor={Roman Czapla, Olaf Bar},
colorlinks=True,
linkcolor=darkgray,  % color of internal links (change box color with linkbordercolor)
citecolor=BrickRed,  % color of links to bibliography
filecolor=Magenta,   % color of file links
urlcolor=BlueViolet}	%%pdfpagemode=FullScreen}

% diagramy, grafy itp.
\usepackage{tikz}
\usetikzlibrary{positioning}
\usetikzlibrary{arrows}
\usetikzlibrary{arrows.meta}
\usetikzlibrary{chains,fit,shapes,calc}
\tikzset{main node/.style={circle,fill=blue!20,draw,minimum size=1cm,inner sep=0pt}}

% algorytmy
\usepackage[linesnumbered,lined,commentsnumbered]{algorithm2e}
\SetKwFor{ForEach}{for each}{do}{end for}%
\SetKwFor{ForAll}{for all}{do}{end for}%
\newenvironment{myalgorithm}
{\rule{\textwidth}{0.5mm}\\\SetAlCapSty{}\SetAlgoNoEnd\SetAlgoNoLine\begin{algorithm}}{\end{algorithm}\rule{\textwidth}{0.5mm}}


%---------------------
\overfullrule=2mm
\pagestyle{plain}
\textwidth=15cm \textheight=685pt \topmargin=-25pt \linespread{1.3} 
\setlength{\parskip}{0pt}
\setlength\arraycolsep{2pt}
\oddsidemargin = 0.9cm
\evensidemargin =-0.1cm

\captionsetup{width=.95\linewidth, justification=centering}
%---------------------




\newtheorem{tw}{Twierdzenie}[section]
\newtheorem{lem}[tw]{Lemat}
\newtheorem{co}[tw]{Wniosek}
\newtheorem{prop}[tw]{Stwierdzenie}
\theoremstyle{definition}
\newtheorem{ex}{Przykład}
\newtheorem{re}[tw]{Uwaga}
\newtheorem{de}{Definicja}[section]



\newcommand{\bC}{{\mathbb C}}
\newcommand{\bR}{{\mathbb R}}
\newcommand{\bZ}{{\mathbb Z}}
\newcommand{\bQ}{{\mathbb Q}}
\newcommand{\bN}{{\mathbb N}}
\newcommand{\captionT}[1]{\caption{\textsc{\footnotesize{#1}}}}
\renewcommand\figurename{Rys.}

\numberwithin{equation}{section}
\renewcommand{\thefootnote}{\arabic{footnote})}
%\renewcommand{\thefootnote}{\alph{footnote})}



\begin{document}

% --------------------------------------------
% Strona tytułowa
% --------------------------------------------

\thispagestyle{empty}
\begin{titlepage}
\begin{center}\Large
Uniwersytet Komisji Edukacji Narodowej w Krakowie\\
\large
Instytut Bezpieczeństwa i Informatyki\\
\vskip 10pt
\end{center}
\begin{center}
\centering \includegraphics[width=1.0\columnwidth]{images/logo.png }
\end{center}

\begin{center}
 {\bf \fontsize{14pt}{14pt}\selectfont PROJEKT INŻYNIERSKI \\ SPRAWOZDANIE}
\end{center}
\vskip 5pt
\begin{center}
 {\bf \fontsize{22pt}{22pt}\selectfont System WasteFree Cloud do efektywnego zarządzania odpadami}
\end{center}

\begin{center}
 {\fontsize{12pt}{12pt}\selectfont wykonany przez: }
\end{center}
\begin{center}
 {\bf\fontsize{16pt}{16pt}\selectfont Dawid Plata}\\
 {\fontsize{12pt}{12pt}\selectfont Nr albumu: 165303 \\\&\\}
 {\bf\fontsize{16pt}{16pt}\selectfont Michał Szymbara}\\
 {\fontsize{12pt}{12pt}\selectfont Nr albumu: 165297\\\&\\}
 {\bf\fontsize{16pt}{16pt}\selectfont Michał Wechmann}\\
 {\fontsize{12pt}{12pt}\selectfont Nr albumu: 165298}
\end{center}
\begin{center}
 {\fontsize{12pt}{12pt}\selectfont pod opieką:}\\
 {\bf\fontsize{12pt}{12pt}\selectfont dr Wojciech Gwizdała}\\
 {\bf\fontsize{12pt}{12pt}\selectfont mgr Wojciech Baran}
\end{center}

%\mbox{}
\vspace*{\fill}
%\vskip 50pt
\begin{center}
\large
Kraków \the\year\\
(ostatnia aktualizacja: \DTMcurrenttime,\;\today)
\end{center}
\end{titlepage}
\setcounter{page}{0} 
\newpage\null\thispagestyle{empty}
%\setcounter{page}{0} 
%\newpage
%\thispagestyle{empty}

\tableofcontents


\newpage

\section{Opis projektu}
Celem projektu było stworzenie systemu informatycznego \textbf{WasteFree Cloud}, wspierającego koordynację odbioru odpadów w modelu grupowym. Aplikacja umożliwia użytkownikom łączenie się w grupy adresowe, wspólne zamawianie kontenerów na odpady, automatyczne rozliczanie kosztów poprzez wbudowany portfel oraz komunikację w czasie rzeczywistym.

System został zrealizowany w architekturze klient-serwer. Warstwa backendowa (API) została napisana w technologii \textbf{.NET 9} z wykorzystaniem \textbf{Minimal API} oraz \textbf{Entity Framework Core}, co zapewnia wysoką wydajność i skalowalność. Komunikacja w czasie rzeczywistym (powiadomienia, czat) oparta jest o bibliotekę \textbf{SignalR}. Warstwa frontendowa (UI) to aplikacja typu SPA (Single Page Application) stworzona w frameworku \textbf{Angular 17}, wykorzystująca nowoczesne mechanizmy takie jak \textbf{Signals} oraz bibliotekę \textbf{RxJS}. Całość rozwiązania jest skonteneryzowana przy użyciu \textbf{Docker} i wdrażana automatycznie za pomocą potoków \textbf{GitHub Actions} na serwer VPS (Mikr.us). Dane plikowe przechowywane są w chmurze \textbf{Azure Blob Storage}, a usługi zewnętrzne obejmują m.in. geokodowanie (Nominatim) i wysyłkę e-maili (SMTP).

\section{Opis zrealizowanych pracy}
W ramach realizacji projektu wykonano następujące prace:
\begin{itemize}
    \item \textbf{Analiza i Projekt}: Opracowanie wymagań funkcjonalnych i niefunkcjonalnych, zaprojektowanie architektury systemu (Clean Architecture), modelu domenowego oraz schematu bazy danych.
    \item \textbf{Backend (.NET)}: Implementacja warstw Domain, Application, Infrastructure i API. Stworzenie logiki biznesowej (mediator, CQRS), mechanizmów autoryzacji (JWT), obsługi zadań w tle (TickerQ) oraz integracji z zewnętrznymi serwisami.
    \item \textbf{Frontend (Angular)}: Implementacja interfejsu użytkownika z podziałem na moduły (Auth, Portal, GarbageAdmin). Opracowanie logiki klienta, serwisów HTTP, interceptorów (Auth, Error, I18n) oraz komponentów interaktywnych.
    \item \textbf{Infrastruktura i DevOps}: Konfiguracja środowiska Docker (Compose i Dockerfile), wdrożenie CI/CD (GitHub Actions), konfiguracja serwera produkcyjnego oraz usług chmurowych.
    \item \textbf{Weryfikacja i Dokumentacja}: Przeprowadzenie testów manualnych, weryfikacja scenariuszy użycia, opracowanie dokumentacji projektowej i użytkowej.
\end{itemize}
Prace były realizowane przez trzyosobowy zespół, z podziałem odpowiedzialności obejmującym zarówno zadania full-stack, jak i specjalizacje w poszczególnych warstwach systemu.
\section{Wykaz załączników}
\textit{Wykaz załączników i ich krótka charakteryzacja.}

\begin{itemize}
\item {binaria i kody źródłowe;}
\item \textit{dokumentacja projektowa;}
\item \textit{dokumentacja użytkowa;}
\item \textit{dokumentacja instalacyjno-konfiguracyjna systemu;}
\item \textit{film dokumentujący działającą aplikację/system.}
\end{itemize}

\section{Oświadczenia}
\begin{center}
\textbf{Oświadczanie o samodzielności wykonania projektu}
\end{center}

Oświadczamy, że niniejszy projekt jest wynikiem samodzielnych prac prowadzonych przez nas pod nadzorem \textit{dr. Wojciecha Gwizdały i mgr. Wojciecha Barana}, Instytut Bezpieczeństwa i Informatyki, Uniwersytet Komisji Edukacji Narodowej w Krakowie. Oświadczamy również, że żadna część tego projektu i jego części nie została lub nie jest składana w~innym miejscu w celu uzyskania tytułu zawodowego/naukowego lub dyplomu.

\vspace*{20pt}
\noindent
\textbf{Dawid Plata} \\
\textbf{Michał Szymbara} \\
\textbf{Michał Wechmann}\\
\vspace*{20pt}

\pagebreak
\begin{center}
\textbf{Oświadczenie o wkładzie pracy poszczególnych członków zespołu projektowego} 
\end{center}

Oświadczamy, że wkład poszczególnych członków zespołu  w powstanie niniejszego projektu inżynierskiego jest następujący:

\begin{table}[H]
\centering
\begin{tabular}{>{\centering}p{3cm}|>{\centering}p{5cm}|>{\centering}p{5cm}c}
                                                      & imię i nazwisko członka zespołu & wkład poszczególnych członków &\\ \hline
\multirow{3}{3cm}{\centering opracowanie koncepcji i~założeń}&
					Dawid Plata&	20 godz.		&\\ \cline{2-4}&
					Michał Szymbara&	20 godz.		&\\ \cline{2-4}&                                   									Michał Wechmann&	20 godz.		&\\ \hline
\multirow{3}{3cm}{\centering opracowanie metod}&
					Dawid Plata&	20 godz.		&\\ \cline{2-4}&
					Michał Szymbara&	30 godz.		&\\ \cline{2-4}&                                   									Michał Wechmann&	20 godz.		&\\ \hline
\multirow{3}{3cm}{\centering przeprowadzenie badań}&
					Dawid Plata&	20 godz.		&\\ \cline{2-4}&
					Michał Szymbara&	20 godz.		&\\ \cline{2-4}&                                   									Michał Wechmann&	20 godz.		&\\ \hline
\multirow{3}{3cm}{\centering wykonanie części praktycznej}&
					Dawid Plata&	110 godz.		&\\ \cline{2-4}&
					Michał Szymbara&	140 godz.		&\\ \cline{2-4}&                                   									Michał Wechmann&	110 godz.		&\\ \hline
\multirow{3}{3cm}{\centering pozostały wkład (sprecyzować rodzaj)}&                                                                        
					Dawid Plata&	dokumentacja, testy (30 godz.)		&\\ \cline{2-4}&
					Michał Szymbara&	dokumentacja, testy (40 godz.)		&\\ \cline{2-4}&                                   									Michał Wechmann&	dokumentacja, testy (30 godz.)		&
\end{tabular}
\end{table}


\vspace*{20pt}
\noindent
\textbf{Dawid Plata} \\
\textbf{Michał Szymbara} \\
\textbf{Michał Wechmann }\\

\renewcommand\refname{Literatura (jeżeli wymagana)}
\bibliography{references}
\addcontentsline{toc}{section}{Literatura}
% --------------------------------------------------------------------
%%%%%%% odkomentować gdy bibliografia ma być wewnątrz dokumentu
% --------------------------------------------------------------------
%\begin{thebibliography}{11}
%
%\addcontentsline{toc}{section}{Literatura}
%
%\bibitem{ZAN}
%C. Zannoni and P. Pasini, 
%\emph{Advances in the Computer Simulatons of Liquid Crystals}, Kluwer Academic Publishers, 2000.
%
%\end{thebibliography}

\end{document}

